\section{Proof of the main theorem} \label{sec: proof}
We first reduce to the case of degree $4$ algebras. Let $A^\bullet$ be a commutative graded $\R$-algebra equipped with a degree map $\deg_A\colon A^d\to\R$ satisfying the Poincaré duality.  For $x\in A$, let $\ann(x)$ denote the annihilator of $x$. When $x$ is nonzero and homogeneous of degree $k$, the quotient ring $A/\ann(x)$ is graded of degree $d-k$, with an induced degree map
\[
    \deg_{A/\ann(x)}(p(y))=\deg_A(xy)
\]
for all $y\in A^{d-k}$, where $p\colon A\to A/\ann(x)$ is the natural projection. We recall the following descent lemma.
\begin{proposition}[{\cite[Lemma 1.16]{CKS87}} and {\cite[Theorem 2.1.5]{KK87}}, see also {\cite[Theorem 3.2]{Cat08}} and {\cite[Lemma 4.5]{AHL25}}]\label{prop: descent}
    $A/\ann(x)$ satisfies Poincaré duality. If furthermore $A$ has the Kähler package with respect to an open convex cone $K_A$ and $x\in \overline K_A$, then $A/\ann(x)$ satisfies the Kähler package with respect to the image of $K_A$.
\end{proposition}

By repeatedly applying \cref{prop: descent}, we see that for $\ell_5,\ldots,\ell_d\in \overline{K}_A\setminus 0$, whose product is nonzero, we have
\[
    \deg_{A/\ann(\ell_5\cdots\ell_d)}(x)=\deg_{A}(x \ell_5\cdots\ell_d)
\]
and that $A/\ann(\ell_5\cdots\ell_d)$ satisfies the K\"ahler package with respect to the projection of $K_A$. To prove \cref{thm: main} for $A$ with respect to $K_A$, it suffices to prove it for $A/\ann(\ell_5\cdots\ell_d)$ with respect to the image of $K_A$.

For the rest of this section, we assume $A=\bigoplus_{i=0}^4 A^i$ which has the K\"ahler package with respect to degree map $\deg$ and open convex cone $K\subseteq A^1$. For any $\alpha,\beta\in K$, we have the Lorentzian bilinear form on $A^1$
\[
    T_{\alpha\beta}(x,y)\coloneqq \deg(\alpha\beta xy).
\]
For any $h\in K$, we can further define the inner product
\[
    G_{\alpha\beta,h}(x,y)\coloneqq G_{T_{\alpha\beta},h}(x,y)=\frac{2 \deg(\alpha\beta h x)\deg(\alpha\beta h y)}{\deg(\alpha\beta h^2)}-\deg(\alpha\beta x y).
\]
We partially order $A^1$ by saying $\alpha \prec \beta$  (or $\alpha \preceq \beta$) when $\beta-\alpha\in K$ (or $\beta-\alpha\in \overline{K}$). Then we have the following interpretation of \cref{lem: estimate of norms}.
\begin{lemma} \label{lem: estimate of norms revisited}
    Let $\alpha,\beta,\gamma\in K$ such that $\alpha\preceq\beta$. Then for every $x\in A^1$,
    \[
        G_{\alpha \gamma,\beta}(x,x)\leq 5G_{\beta \gamma,\beta}(x,x).
    \]
\end{lemma}
\begin{proof}
    In the case where $\alpha\prec \beta$, let $R(x,y)=T_{\alpha \gamma}(x,y), S(x,y)=T_{(\beta-\alpha)\gamma}(x,y)$ and $T(x,y)=T_{\beta \gamma}(x,y)$. It's clear that
    \[
        T=R+S,\ R(\beta,x)=T(\alpha,x),\ S(\beta,x)=T(\beta-\alpha,x)
    \]
    and that
    \[
        R(\beta,\beta),\ S(\beta,\beta),\ T(\alpha,\alpha),\ T(\beta-\alpha,\beta-\alpha)>0.
    \]
    Therefore, we can apply \cref{lem: estimate of norms} and obtain the desired inequality.
    In general, we have inequality
    \[
        G_{(1-\epsilon)\alpha \gamma,\beta}(x,x)\leq 5G_{\beta \gamma,\beta}(x,x),
    \]
    for $0<\epsilon<1$. Taking the limit $\epsilon\to 0$ supplies the desired inequality.
\end{proof}
By Poincar\'e duality, for any $h\in K$, we are allowed to define the dual norm of $G_{h^2}$ on $A^3$:
\[
    \|z\|_{h}^2=\sup_{0\neq x\in A^1}\frac{\deg(zx)^2}{G_{h^2,h}(x,x)}.
\]
Since the inner product $G_{h^2,h}$ is induced by the isomorphism between $A^1$ and $A^3\cong (A^1)^\vee $ given by the composition of multiplication by $h^2$ and $\theta_{T_{h^2},h}$, we see that
\[
    \|h^2x\|_h^2=G_{h^2,h}(\theta_{T_{h^2},h}(x),\theta_{T_{h^2},h}(x))=G_{h^2,h}(x,x),
\]
for any $x\in A^1$. We also have the following estimate of the norm $\|\cdot\|_{h}$.
\begin{lemma}\label{lem: estimate of external norm}
    Let $h,\alpha,\beta\in K$ such that $\alpha,\beta\preceq h$. Then we have
    \[
        \|\alpha\beta x\|_h^2\leq 25 G_{\alpha\beta,h}(x,x)
    \]
    for every $x\in A^1$.
\end{lemma}

\begin{proof}
    By applying \cref{lem: estimate of norms revisited} twice, we get
    \[
        G_{\alpha\beta,h}(y,y)\leq 5G_{h\beta,h}(y,y)\leq 25 G_{h^2,h}(y,y),
    \]
    and by \cref{lem: cauchy-schwarz},
    \[
        \deg(\alpha\beta x y)^2=T_{\alpha\beta}(x,y)^2\leq G_{\alpha\beta,h}(x,x)G_{\alpha\beta,h}(y,y)\leq 25 G_{\alpha\beta,h}(x,x)G_{h^2,h}(y,y).
    \]
    After dividing by $G_{h^2,h}(y,y)$ and taking the supremum over $y$, we obtain the desired inequality.
\end{proof}

Now we are ready to prove the main theorem using the Lefschetz decomposition and the above estimates.
\begin{proof}[Proof of \cref{thm: main}]
Let $R=\R[X_1,\ldots,X_n]$ and $P, H, V_i, A_i, B_i$ be as in the setup of \cref{thm: main}. Fix a graded ring homomorphism $R\to A$ that sends $K_R$ to $K$, and denote the images of $P,H,V_i,A_i,B_i$ by $p,h,v_i,\alpha_i,\beta_i$ respectively.

We first prove the theorem when $0\prec \alpha_i,\beta_i\preceq h$. According to the Lefschetz decomposition theorem, we can write $p$ uniquely as
\[
    p=h^2p_0+hp_1+p_2, \quad p_i\in A^i
\]
such that
\[
    h^3p_1=hp_2=0.
\]
By orthogonality of the Lefschetz decomposition, we have
\[
    \deg(p^2)=\deg(h^4p_0^2)+\deg(h^2p_1^2)+\deg(p_2^2)=p_0^2\deg(h^4)-G_{h^2,h}(p_1,p_1)+\deg(p_2^2).
\]
Meanwhile,
\[
    \|hp\|_h^2=\|h^2(hp_0+p_1)\|_h^2=G_{h^2,h}(hp_0+p_1,hp_0+p_1)=p_0^2\deg(h^4)+G_{h^2,h}(p_1,p_1).
\]
By adding the above equations and applying the Hodge-Riemann relation in degree $2$, we get
\[
    \deg(p^2)+\|hp\|_h^2=2p_0^2\deg(h^4)+\deg(p_2^2)\geq 2p_0^2\deg(h^4)=\frac{2\deg(ph^2)^2}{\deg(h^4)}.\tag{2}\label{eq: external inequality}
\]
Now apply the triangle inequality and \cref{lem: estimate of external norm} on the image of equality \eqref{eq: referencing}, we have
\[
    \|hp\|_h\leq \sum_{i=1}^r\|\alpha_i\beta_iv_i\|_h\leq\sum_{i=1}^r 5\sqrt {G_{\alpha_i\beta_i,h}(v_i,v_i)}=5\sum_{i=1}^r\sqrt{\frac{2\deg(\alpha_i\beta_ihv_i)^2}{\deg(\alpha_i\beta_ih^2)}-\deg(\alpha_i\beta_iv_i^2)}.
\]
By substituting the estimate above into inequality \eqref{eq: external inequality} we obtain the desired inequality.

In general, when $0\preceq\alpha_i,\beta_i\preceq h$ such that $\deg(\alpha_i\beta_ih^2)>0$, we can use a limit argument as follows: we can define perturbations of the chosen morphism $\varphi\colon R\to A$ as
\[
    \varphi_{\epsilon}\colon R\to A,\;X_i\mapsto\varphi(X_i)+\epsilon \ell
\]
where $\ell\in  K$ is any fixed ample class. Then for $0<\epsilon\ll 1$, $\varphi_\epsilon$ sends $\overline{K}_R\setminus0\to K$. In particular, if we denote the images of $P,H,V_i,A_i,B_i$ by $p_\epsilon,h_\epsilon,v_{i,\epsilon},\alpha_{i,\epsilon},\beta_{i,\epsilon}$, we have shown that
\[
   \deg(p_\epsilon^2)+25\left(\sum_{i=1}^r\sqrt{\frac{2\deg(\alpha_{i,\epsilon}\beta_{i,\epsilon}h_{\epsilon}v_{i,\epsilon})^2}{\deg(\alpha_{i,\epsilon}\beta_{i,\epsilon}h_{\epsilon}^2)}-\deg(\alpha_{i,\epsilon}\beta_{i,\epsilon}v_{i,\epsilon}^2)}\right)^2\geq \frac{2\deg(p_\epsilon h_\epsilon^2)^2}{\deg(h_\epsilon^4)},
\]
note that the assumption on the denominators ensures $A_i,B_i,H\neq 0$. Since both sides are continuous in $\epsilon$, after taking the limit of $\epsilon\to 0$, we obtain the desired inequality. Note that \[\deg(h^4)=\deg(\alpha_i\beta_ih^2)+\deg((h-\alpha_i)\beta_ih^2)+\deg(\alpha_i(h-\beta_i)h^2)+\deg((h-\alpha_i)(h-\beta_i)h^2)>0,\]
so the right-hand side has a well-defined limit.
\end{proof}

\begin{remark}
\begin{enumerate}[leftmargin=*]
    \item When $r=1$ and $A_1=B_1=H$, the theorem degenerates to
    \[\deg(h^2v_1^2)+25\left(\frac{2\deg(h^3v_1)^2}{\deg(h^4)}-\deg(h^2v_1^2)\right)\geq\frac{2\deg(h^3v_1)^2}{\deg(h^4)},\]
    which can be easily deduced from the AF-type inequality.
    \item 
    The estimate of the left-hand side of \eqref{eq: external inequality} is necessary. If $R\to A$ is not onto, the preimage of $hp$ under the Lefschetz isomorphism might not be contained in the image of $R$, and hence the value of $\|hp\|_h$ could be unknown. The current estimate is far from being sharp, but we'll see it is enough for the purpose of section \ref{sec: ex}.
    \item If we allow $A$ and $p$ to have higher degrees, one can still utilize the Lefschetz decomposition and repeatedly construct norms on $A^i$, and it's not hard to write down similar inequalities like \eqref{eq: external inequality}. However, obtaining analogous norm estimates like \Cref{lem: estimate of external norm} remains an open problem.
\end{enumerate}
\end{remark}